\subsection{与滤过相容的同态}

	设 \(G, G'\) 为两个交换群（用加法记号），\((G_n)\) 为 \(G\) 上的滤过，\((G'_n)\) 为 \(G'\) 上的滤过；同态 \(h: G \to G'\) 称为与 \(G\) 和 \(G'\) 上的滤过相容，如果 \(h(G_n) \subset G'_n\) 对所有 \(n \in \mathbf{Z}\) 成立。复合同态 \(G_n \xrightarrow{h|G_n} G'_n \longrightarrow G'_n / G_{n+1}'\) 在 \(G_{n+1}\) 上为零，因而取商后定义出同态 \(h_n: G_n / G_{n+1} \to G'_n / G_{n+1}'\)；于是存在唯一的加法群同态 \(\text{gr}(h): \text{gr}(G) \to \text{gr}(G')\)，使得对所有 \(n \in \mathbf{Z}\)，\(\text{gr}(h)\) 与 \(h_n\) 在 \(\text{gr}_n(G) = G_n / G_{n+1}\) 上重合。\(\text{gr}(h)\) 称为与 \(h\) 相伴的分次群同态。若 \(G''\) 是第三个滤过群，且 \(h': G' \to G''\) 是与滤过相容的同态，则 \(h' \circ h\) 也是与滤过相容的同态，并且

\[
\operatorname{gr} (h ^ {\prime} \circ h) = \operatorname{gr} (h ^ {\prime}) \circ \operatorname{gr} (h).\tag{12}
\]

命题 2. 设 G 为滤过交换群，H 为 G 的子群；给 H 赋予诱导滤过，给 G/H 赋予商滤过。若 j: H → G 是典范嵌入，p: G → G/H 是典范满射，则 j 和 p 与滤过相容，且序列

\[
0 \longrightarrow \operatorname{gr} (\mathrm{H}) \xrightarrow {\operatorname{gr} (j)} \operatorname{gr} (\mathrm{G}) \xrightarrow {\operatorname{gr} (p)} \operatorname{gr} (\mathrm{G} / \mathrm{H}) \longrightarrow 0\tag{13}
\]

是正合的。

第一个断言显然；若 \((\mathbf{G}_n)\) 是 \(\mathbf{G}\) 上的滤过，则

\[
\left(\mathrm{H} \cap \mathrm{G} _ {n}\right) \cap \mathrm{G} _ {n + 1} = \mathrm{H} \cap \mathrm{G} _ {n + 1}
\]

从而 \(\mathrm{gr}(j)\) 是单射；此外，典范映射

\[
\mathrm{G} _ {n} \rightarrow (\mathrm{H} + \mathrm{G} _ {n}) / \mathrm{H}
\]

是满射，从而 \(\operatorname{gr}(p)\) 也是满射，并由 (12) 有 \(\operatorname{gr}(p) \circ \operatorname{gr}(j) = 0\)。最后，设 \(\xi \in \mathbf{G}_n / \mathbf{G}_{n+1}\) 属于 \(\operatorname{gr}(p)\) 的核；则存在 \(x \in \xi\) 使得 \(x \in \mathbf{H} + \mathbf{G}_{n+1}\)；但由于 \(\mathbf{G}_{n+1} \subset \mathbf{G}_n\)，

\[
\mathrm{G} _ {n} \cap (\mathrm{H} + \mathrm{G} _ {n + 1}) = (\mathrm{H} \cap \mathrm{G} _ {n}) + \mathrm{G} _ {n + 1}
\]

	因此 \(x = y + z\)，其中 \(y \in \mathbf{H} \cap \mathbf{G}_n\)，\(z \in \mathbf{G}_{n+1}\)；这说明 \(\xi\) 是 \(\mathbf{G}_{n+1}\) 模下 \(j(y)\) 的类，换言之，它属于 \(\operatorname{gr}(\mathbf{H})\) 在 \(\operatorname{gr}(j)\) 下的像。

注意，对于滤过交换群的正合序列 \(0 \to G' \xrightarrow{u} G \xrightarrow{v} G'' \to 0\)，若 \(u\) 和 \(v\) 与滤过相容，则序列 \(0 \longrightarrow \text{gr}(G') \xrightarrow{\text{gr}(u)} \text{gr}(G) \xrightarrow{\text{gr}(v)} \text{gr}(G'') \longrightarrow 0\) 未必正合（习题 4）。

现在若 A 和 B 是两个滤过环，且 \(h \colon \mathbf{A} \to \mathbf{B}\) 是与滤过相容的环同态，则立即验证可知，分次群同态 \(\mathrm{gr}(h) \colon \mathrm{gr}(\mathbf{A}) \to \mathrm{gr}(\mathbf{B})\) 也是环同态。特别地，若 \(\mathbf{A}'\) 是带诱导滤过的 A 的子环，则 \(\mathrm{gr}(\mathbf{A}')\) 可典范地识别为 \(\mathrm{gr}(\mathbf{A})\) 的分次子环（命题 2）；若 \(\mathfrak{b}\) 是 A 的双边理想，且给 \(\mathbf{A} / \mathfrak{b}\) 赋予商滤过，则 \(\mathrm{gr}(\mathbf{A} / \mathfrak{b})\) 可典范地识别为商分次环 \(\mathrm{gr}(\mathbf{A}) / \mathrm{gr}(\mathfrak{b})\)（命题 2）。

最后，设 A 为滤过环，E、F 为两个滤过 A-模，且 \(u: \mathbf{E} \to \mathbf{F}\) 是与滤过相容的同态。立即可见，\(\mathrm{gr}(u): \mathrm{gr}(\mathrm{E}) \to \mathrm{gr}(\mathrm{F})\) 是 \(\mathrm{gr}(\mathrm{A})\)-线性映射，因而是分次 \(\mathrm{gr}(\mathrm{A})\)-模的一个次数为 0 的齐次同态。此外，若 \(u': \mathbf{E} \to \mathbf{F}\) 是另一个与滤过相容的 A-同态，则 \(u + u'\) 也与滤过相容，并且

\[
\operatorname{gr} (u + u ^ {\prime}) = \operatorname{gr} (u) + \operatorname{gr} (u ^ {\prime}).\tag{14}
\]

\textbf{注}\par

\begin{enumerate}
\def\labelenumi{(\arabic{enumi})}
\item
  显然，与滤过相容的滤过环同态（分别地，给定滤过环 A 上的滤过模同态）可以作为滤过环结构（分别地，滤过 A-模结构）的态射（《集合论》，第 IV 章，§ 2，第 1 节）。
\item
  设 E 和 F 是滤过环 A 上的两个模，并给它们赋予由 A 上的滤过 \((\mathbf{A}_{n})\) 导出的滤过（第 1 节，例 2）。则每个 A-线性映射 \(u: E \rightarrow F\) 都与滤过相容，因为

\[
u (\mathrm{A} _ {n} \mathrm{E}) = \mathrm{A} _ {n} u (\mathrm{E}) \subset \mathrm{A} _ {n} \mathrm{F}.
\]

\item
	  注意，一个与滤过相容的滤过 A-模同态 \(u: \mathbf{E} \to \mathbf{F}\) 可能满足 \(\operatorname{gr}(u) = 0\) 而自身不为零；例如，加法群的自同态 \(x \mapsto nx\) 作用在 \(\mathbf{Z}\) 上，而其带有 \((n)\)-进滤过（对任意整数 \(n > 1\)）。因此，关系 \(\operatorname{gr}(u) = \operatorname{gr}(v)\)（其中 \(u, v\) 是从 \(\mathbf{E}\) 到 \(\mathbf{F}\) 的两个与滤过相容的同态）并不必然推出 \(u = v\)。
\item
  本节开头的定义立即推广到两个不必交换的群 G、G'：它们分别由子群 Gn、Gn' 滤过，且 Gn+1（分别 Gn'+1）是 Gn（分别 Gn'）的正规子群。在对 Gn 作相同假设并假定 H 在 G 中不变时，命题 2 也成立，证明除记号外不变。
\end{enumerate}

